%% ===================================================================
%% universal-bridge-attestation-soundness.tex
%%
%% The B-Chain (transport/relay) o M-Chain (M-of-N threshold
%% attestation) composition, and why the M/T/F separation of concerns
%% is a load-bearing security boundary (not cosmetic).
%%
%% Companion to:
%%   - quasar-cert-soundness.tex  (source-chain finality; the
%%     No-double-finalization theorem is the finality-soundness input)
%%   - LP-6000 (B-Chain / bridgevm), LP-7100 (M-Chain / MPC ceremonies),
%%     LP-8200 (F-Chain / FHE), LP-7050 (chain-topology supersession).
%% ===================================================================

\documentclass[11pt]{article}
\usepackage{amsmath,amssymb,amsthm}
\usepackage{booktabs}
\usepackage{enumitem}
\usepackage{hyperref}
\usepackage[margin=1in]{geometry}
\hypersetup{colorlinks=true,linkcolor=black,citecolor=blue,urlcolor=blue}

\newtheorem{theorem}{Theorem}
\newtheorem{lemma}[theorem]{Lemma}
\newtheorem{corollary}[theorem]{Corollary}
\newtheorem{definition}[theorem]{Definition}
\newtheorem{remark}[theorem]{Remark}

\title{Universal Bridge Attestation Soundness\\
{\large B-Chain $\circ$ M-Chain, and the M/T/F separation as a
security boundary}}
\author{Lux Industries, Inc.\\ \texttt{research@lux.network}}
\date{2026-07-02}

\begin{document}
\maketitle

\begin{abstract}
The Lux universal bridge honours a cross-chain message only when an
$M$-of-$N$ \emph{threshold attestation}, produced by the M-Chain over a
message that is \emph{final on the source chain}, is presented on the
destination chain. We formalise this composition and prove
\emph{bridge attestation soundness}: forging acceptance of a message
that was never source-finalised reduces to either a threshold-signature
forgery or a source-chain finality violation. Both are negligible ---
the first by the EUF-CMA security of the M-Chain signing scheme
(CGGMP21 / FROST), the second by the QuasarCert
\emph{No double-finalization} theorem
(\texttt{quasar-cert-soundness.tex}~\S``Finality Admission''). We then
show that this reduction \emph{depends on} the M/T/F separation of
concerns (LP-7050, LP-7100, LP-8200): braiding the M-Chain signing
shares with F-Chain FHE compute in one VM shares a fault domain between
the bridge-custody key and the encrypted-compute path, silently
weakening the unforgeability assumption the whole reduction rests on.
Separation is therefore a security boundary, not an aesthetic one.
\end{abstract}

\section{Chains and concerns}
\label{sec:chains}

\begin{definition}[Concern-owning chains]
\label{def:mtf}
Per LP-134 / LP-7050 the roster is $P,C,X,Q,Z,A,B,M,F$, and each of the
former T-Chain concerns has exactly one owner:
\begin{description}[nosep,leftmargin=2em]
  \item[M-Chain (MVM, LP-7100):] runs $M$-of-$N$ MPC signing ceremonies
        (CGGMP21 threshold ECDSA, FROST threshold Schnorr, Pulsar-general)
        for bridge custody of external-chain wallets. It \emph{attests};
        it does not decide finality and does not compute on ciphertexts.
  \item[F-Chain (FVM, LP-8200):] runs FHE (TFHE bootstrap-key DKG,
        encrypted-EVM execution, no-reconstruct threshold decryption). It
        \emph{computes on encrypted data}; it holds no bridge-custody
        signing key.
  \item[B-Chain (\texttt{bridgevm}, LP-6000):] transports and relays
        cross-chain messages (teleport). It \emph{carries} attestations;
        it mints no signature and finalises no source block.
  \item[Source chain $S$ (any Lux L1/L2):] produces QuasarCert finality
        over its own blocks (\S\ref{sec:source}).
\end{description}
The M-Chain and F-Chain \emph{share} the \texttt{thresholdvm} substrate
library but are distinct chains; \S\ref{sec:separation} shows why the
distinctness is load-bearing.
\end{definition}

\section{Source-chain finality (the input)}
\label{sec:source}

\begin{definition}[Source finality]
\label{def:source-final}
For an event slot $e$ on source chain $S$ (a height, or a
height-and-log-index), let $\mathrm{Final}_S(e, m) = 1$ iff message $m$
is committed at $e$ by a verified QuasarCert over the canonical
execution commitment at $e$'s height.
\end{definition}

\begin{lemma}[Finality soundness of $S$]
\label{lem:source-sound}
Below $S/3$ Byzantine stake on $S$, for every event slot $e$ there is at
most one $m$ with $\mathrm{Final}_S(e,m)=1$.
\end{lemma}
\begin{proof}
Immediate from the QuasarCert \emph{No double-finalization} theorem
(\texttt{quasar-cert-soundness.tex}~\S``Finality Admission''): two
certificates at one height over different canonical commitments require
an honest validator to have signed both, which the durable
decided-height sign gate forbids. A single canonical commitment fixes a
single $m$ at $e$.
\end{proof}

\section{The M-Chain attestation and the B-Chain relay}
\label{sec:attest}

\begin{definition}[Threshold attestation scheme]
\label{def:ts}
The M-Chain committee of size $N$ holds an $M$-of-$N$ threshold key with
aggregate public key $\mathsf{apk}$ (established by a DKG; no party holds
the secret). $\mathsf{TS}.\mathsf{Sign}$ is the interactive ceremony
requiring $\ge M$ honest participants; $\mathsf{TS}.\mathsf{Verify}(
\mathsf{apk}, x, \sigma)\to\{0,1\}$. $\mathsf{TS}$ is EUF-CMA: an
adversary controlling $<M$ committee members and making chosen-message
ceremony queries forges a signature on a fresh $x$ with advantage at
most $\epsilon_{\mathsf{TS}}$ (CGGMP21 under enhanced-existential
unforgeability of ECDSA; FROST under the OMDL / AGM analyses).
\end{definition}

\begin{definition}[Attested bridge message and acceptance]
\label{def:accept}
The M-Chain emits an attestation only after checking source finality:
\[
  x = (\mathsf{srcID}, e, m, \mathsf{finalityProof}),\quad
  \sigma \leftarrow \mathsf{TS}.\mathsf{Sign}(x)
  \ \text{iff}\ \mathrm{Final}_S(e,m)=1 .
\]
B-Chain relays $(x,\sigma)$ to destination $D$. $D$ \textbf{accepts} $m$
for $e$ iff
\[
  \mathsf{TS}.\mathsf{Verify}(\mathsf{apk}, x, \sigma)=1
  \ \wedge\ x.m = m
  \ \wedge\ e \notin \mathsf{consumed}_D,
\]
and then records $e \in \mathsf{consumed}_D$ (single-spend / replay
guard).
\end{definition}

\section{Soundness}
\label{sec:soundness}

\begin{theorem}[Bridge attestation soundness]
\label{thm:bridge}
Assume (i) fewer than $M$ of the $N$ M-Chain committee members are
Byzantine, (ii) $S$ has $<S/3$ Byzantine stake, and (iii)
$\mathsf{TS}$ is EUF-CMA with advantage $\epsilon_{\mathsf{TS}}$. Then
the probability that destination $D$ accepts, for some event slot $e$,
a message $m^\star$ with $\mathrm{Final}_S(e,m^\star)=0$ is at most
$\epsilon_{\mathsf{TS}}$. Moreover $D$ accepts at most one $m$ per $e$,
and that $m$ is the unique source-finalised message for $e$.
\end{theorem}
\begin{proof}
Suppose $D$ accepts $m^\star$ for $e$. By
Definition~\ref{def:accept} there is a $\sigma^\star$ with
$\mathsf{TS}.\mathsf{Verify}(\mathsf{apk}, x^\star, \sigma^\star)=1$ and
$x^\star.m = m^\star$. Consider the message $x^\star = (\mathsf{srcID},
e, m^\star, \cdot)$.

\emph{Case A: $x^\star$ was never signed by the honest ceremony.} Then
$\sigma^\star$ is a valid signature on a message the honest committee
never produced. By assumption (i) the adversary controls $<M$ members,
so producing $\sigma^\star$ is exactly an EUF-CMA forgery against
$\mathsf{TS}$; this occurs with probability $\le \epsilon_{\mathsf{TS}}$.

\emph{Case B: $x^\star$ was signed by the honest ceremony.} By
Definition~\ref{def:accept} the ceremony signs $x^\star$ only when
$\mathrm{Final}_S(e, m^\star)=1$. Hence $m^\star$ \emph{is}
source-finalised, contradicting the hypothesis
$\mathrm{Final}_S(e,m^\star)=0$.

So a wrongful acceptance implies Case~A, of probability
$\le\epsilon_{\mathsf{TS}}$. For the uniqueness clause: any accepted $m$
falls in Case~B, so $\mathrm{Final}_S(e,m)=1$; by
Lemma~\ref{lem:source-sound} at most one such $m$ exists, and the
$\mathsf{consumed}_D$ guard admits $e$ only once, so $D$ accepts exactly
that unique $m$ and no replay. $\qed$
\end{proof}

\begin{remark}[Two disjoint trust hypotheses]
\label{rem:two-hyp}
Bridge soundness is the \emph{conjunction} of two independent
assumptions living on two different chains: threshold honesty on the
M-Chain (hypothesis i) and stake honesty on the source chain
(hypothesis ii). Neither implies the other, and each is verified in its
own domain --- exactly the ``one concern, one chain'' discipline of
Definition~\ref{def:mtf}. A single actor cannot forge a bridge
acceptance without either compromising $M$ M-Chain signers or breaking
source finality.
\end{remark}

\section{Why separation of concerns is load-bearing}
\label{sec:separation}

Theorem~\ref{thm:bridge} rests on hypothesis (iii): $\mathsf{TS}$ is
EUF-CMA. That, in turn, presumes the threshold \emph{signing shares} are
confined to the M-Chain's fault domain --- an adversary that can read or
fault-inject a share does not need $M$ colluding members.

\begin{definition}[Fault domain]
\label{def:fault-domain}
A \emph{fault domain} is a unit of code + memory + accelerator state
within which a fault (memory-disclosure bug, side channel, malicious
GPU kernel, shared-address-space corruption) can move secret material.
Two subsystems in the same VM/process share a fault domain.
\end{definition}

\begin{lemma}[Braiding enlarges the trusted surface]
\label{lem:braid}
Let $\mathsf{TS}$'s signing shares and an FHE evaluation engine reside
in one VM (the overloaded \texttt{thresholdvm}: MPC signing $+$
\texttt{vms/thresholdvm/fhe} in one address space). Then the
unforgeability advantage against $\mathsf{TS}$ is bounded not by
$\epsilon_{\mathsf{TS}}$ alone but by
$\epsilon_{\mathsf{TS}} + \epsilon_{\mathrm{FHE\text{-}leak}}$, where
$\epsilon_{\mathrm{FHE\text{-}leak}}$ is the probability that a fault in
the FHE path (a bug in an encrypted-EVM opcode, a malicious bootstrap
kernel, a decryption side channel) discloses or corrupts a signing
share.
\end{lemma}
\begin{proof}
An adversary now has two ways to obtain a forgery: the $M$-collusion
path (bounded by $\epsilon_{\mathsf{TS}}$) \emph{or} exercising the FHE
path to exfiltrate/fault a share and sign directly (bounded by
$\epsilon_{\mathrm{FHE\text{-}leak}}$, since both live in one fault
domain by Definition~\ref{def:fault-domain}). The advantage is the
union bound $\epsilon_{\mathsf{TS}} + \epsilon_{\mathrm{FHE\text{-}leak}}$.
The FHE surface is large (a full encrypted-EVM opcode set plus GPU
bootstrap kernels), so $\epsilon_{\mathrm{FHE\text{-}leak}}$ is neither
zero nor small a priori.
\end{proof}

\begin{corollary}[M/F separation is a security boundary]
\label{cor:separation}
Placing MPC signing (M-Chain, LP-7100) and FHE compute (F-Chain,
LP-8200) in \emph{distinct} chains with \emph{disjoint} fault domains
restores hypothesis (iii) to $\epsilon_{\mathsf{TS}}$ alone, and thus
restores Theorem~\ref{thm:bridge} to its stated bound. The separation is
therefore not cosmetic: it removes the $\epsilon_{\mathrm{FHE\text{-}leak}}$
term from bridge soundness. The two chains may share the
\texttt{thresholdvm} substrate \emph{library} and even
\texttt{constants.ThresholdVMID} (disambiguated by
\texttt{CreateChainTx.BlockchainName}, LP-134) without sharing a fault
domain, provided the running instances do not co-reside signing shares
with FHE evaluation state.
\end{corollary}

\begin{remark}[Implementation status]
\label{rem:impl-status}
As of 2026-07-02 the node has not yet performed the split: the genesis
registry (\texttt{node/genesis/builder/registry.go}) lists a single
overloaded \texttt{T-Chain} entry and no M/F rows, and the
\texttt{thresholdvm} package still braids MPC signing with
\texttt{vms/thresholdvm/fhe}. Corollary~\ref{cor:separation} is thus the
\emph{target} security posture; LP-7050~\S8 tracks the migration. Until
the split lands, bridge soundness carries the
$\epsilon_{\mathrm{FHE\text{-}leak}}$ term of Lemma~\ref{lem:braid} and
should be treated accordingly (keep the FHE surface off any node that
holds live M-Chain custody shares).
\end{remark}

\section{Conclusion}
The universal bridge is sound exactly when two disjoint honesty
hypotheses hold on two different chains, and its soundness proof
\emph{requires} the M/T/F separation to keep the threshold-signature
assumption clean. Separation of concerns here is a theorem hypothesis,
not a style preference.

\end{document}
